\documentclass[11pt,a4paper]{article}
% Spolecny styl pro ucitelske pripravy MAT (Work/MAT)
% Verze 1.1, 2026-10-23
\NeedsTeXFormat{LaTeX2e}
\ProvidesPackage{priprava}[2026/10/23 v1.1 Sablona ucitelskych priprav MAT]

% Zakladni typografie a cestina
\RequirePackage[T1]{fontenc}
\RequirePackage[utf8]{inputenc}
\RequirePackage[czech]{babel}
\RequirePackage{lmodern}
\RequirePackage{microtype}
\RequirePackage{amsmath,amssymb}
\RequirePackage{xcolor}
\RequirePackage{enumitem}
\RequirePackage[most]{tcolorbox}
\RequirePackage{tabularx,booktabs,array}
\RequirePackage{fancyhdr}
\RequirePackage{needspace}
\RequirePackage{geometry}
\RequirePackage[hidelinks]{hyperref}

\geometry{a4paper,top=17mm,bottom=18mm,left=23mm,right=23mm}
\definecolor{notegray}{gray}{0.39}
\setlength{\parindent}{0pt}
\setlength{\parskip}{2.5pt plus 1pt minus 0.5pt}
\setlist[itemize]{topsep=2pt,itemsep=1pt,parsep=0pt,leftmargin=1.45em}
\setlist[enumerate]{topsep=2pt,itemsep=2pt,parsep=0pt,leftmargin=1.65em}
\setcounter{secnumdepth}{0}
\renewcommand{\arraystretch}{1.12}

% Zahlavi: vlevo rocnik | globalni tema, vpravo cislo: nazev hodiny
% Zapati zustava zcela prazdne.
\pagestyle{fancy}
\fancyhf{}
\renewcommand{\headrulewidth}{0.35pt}
\renewcommand{\footrulewidth}{0pt}
\setlength{\headheight}{14pt}
\newcommand{\nastavPripravu}[4]{%
  \fancyhead[L]{\small #1\enspace\textbar\enspace #2}%
  \fancyhead[R]{\small #3: #4}%
}

% Minimalisticke cernobile boxy. Kazda definice zvlast.
\tcbset{commonlesson/.style={enhanced,breakable=false,colback=white,colframe=black,
  colbacktitle=white,coltitle=black,boxrule=0.45pt,
  arc=0.8mm,outer arc=0.8mm,left=3mm,right=3mm,top=1.6mm,bottom=1.6mm,
  toptitle=1mm,bottomtitle=0.7mm,fonttitle=\bfseries,
  attach boxed title to top left={xshift=2.5mm,yshift=-1.15mm},
  boxed title style={colback=white,colframe=white,boxrule=0pt,left=1mm,right=1mm}}}
\newtcolorbox{zadaniBox}[1][]{commonlesson,title={Zadání},#1}
\newtcolorbox{prikladBox}[1][]{commonlesson,title={Příklad},#1}
\newtcolorbox{definiceBox}[1][]{commonlesson,title={Definice},#1}

% Spolecny box pro zadani a jeho reseni, bez mezery mezi castmi.
% Pouziti:
%   \begin{prikladReseniBox}
%     Zadani ulohy.
%     \tcblower
%     \begin{align*} ... \end{align*}
%   \end{prikladReseniBox}
% Spodni cast ma jemny sedy podklad a delici linku; lze v ni sazet
% vice radku matematiky. Obe casti maji jen jeden spolecny ramecek.
\tcbset{solutionlesson/.style={
  commonlesson,
  enhanced jigsaw,
  breakable,
  colbacklower=black!3,
  segmentation style={solid,draw=black!25,line width=0.35pt},
  fontlower=\small\color{notegray},
  before lower={\noindent\textbf{Řešení:}\enspace},
  leftlower=3mm,rightlower=3mm
}}
\newtcolorbox{prikladReseniBox}[1][]{solutionlesson,title={Příklad},#1}
\newtcolorbox{zadaniReseniBox}[1][]{solutionlesson,title={Zadání},#1}

% Pomucky pro autora pripravy: sedive poznamky; \reseni pro starsi pripravy
\newcommand{\poznamka}[1]{\par\begingroup\small\color{notegray}\itshape
  \textnormal{\textbf{Poznámka pro učitele.}}\enspace #1\par\endgroup}
\newcommand{\reseni}[1]{\par\begingroup\small\color{notegray}
  \textbf{Očekávané řešení:}\enspace #1\par\endgroup}
\newcommand{\blok}[1]{\needspace{5\baselineskip}\subsection*{#1}}
\newcommand{\vyrok}[1]{\uv{#1}}

\nastavPripravu{1. ročník}{Algebraické výrazy}{44}{Násobení mnohočlenu mnohočlenem}
\begin{document}

% Zdroje: vlastni ucebnice Algebraicke vyrazy, Nasobeni polynomu;
% Realisticky.cz 010802 Nasobeni mnohoclenu; navaznost na 43.
\poznamka{Jádro 45 minut: připomenutí 5 min, dvojčlen krát dvojčlen 10 min, tabulka 7 min, samostatná práce 15 min a společná kontrola 8 min. Obtížnější výrazy a rezerva jsou volitelné. Algebraické vzorce se budou systematicky probírat až v 46. hodině.}
\blok{Od známého pravidla ke spolehlivému postupu}
V minulé hodině jsme odvodili $(a+b)(c+d)=ac+ad+bc+bd$. Dnes se zaměříme na \textbf{pořádek v zápisu}, znaménka a součiny, v nichž se členy po roznásobení slučují. Cíl není znát nové pravidlo, ale umět použít distributivitu bez vynechaného součinu.
\begin{definiceBox}[title={Pravidlo násobení mnohočlenů}]
\textbf{Součin dvou mnohočlenů} získáme tak, že každý člen prvního mnohočlenu vynásobíme každým členem druhého mnohočlenu a vzniklé součiny sečteme. Nakonec sloučíme stejnorodé členy.
\end{definiceBox}
\begin{prikladReseniBox}[title={Příklad 1: součin dvou dvojčlenů}]
Uprav $(3x-2)(2x+5)$.
\tcblower
\begin{align*}
(3x-2)(2x+5)&=3x(2x+5)-2(2x+5)\\
&=6x^2+15x-4x-10\\
&=6x^2+11x-10.
\end{align*}
\end{prikladReseniBox}
\blok{Jak předejít vynechání členu?}
\begin{prikladReseniBox}[title={Příklad 2: dvojčlen krát trojčlen}]
Uprav $(2x-3)(x^2+x-2)$ roznásobením závorek.
\tcblower
\begin{align*}
(2x-3)(x^2+x-2)
&=2x(x^2+x-2)-3(x^2+x-2)\\
&=2x^3+2x^2-4x-3x^2-3x+6\\
&=2x^3-x^2-7x+6.
\end{align*}
\end{prikladReseniBox}
\poznamka{Požaduj přehledný mezikrok se všemi součiny; nejprve násobíme celou závorku členem $2x$, potom členem $-3$. U dvojčlenu a trojčlenu musí vzniknout šest dílčích součinů, i když se později sloučí.}
\blok{Samostatná práce — základ}
\begin{zadaniReseniBox}[title={S1: od dvojčlenů k trojčlenům}]
\begin{enumerate}[label=\alph*)]
\item $(x+4)(x-3)$;
\item $(2x-5)(x+1)$;
\item $(a-3)(a^2+a-2)$;
\item $(2x+y)(x-2y)$;
\item $(x+2)(x-3)-(x-1)(x+4)$.
\end{enumerate}
\tcblower
\begin{enumerate}[label=\alph*)]
\item $x^2+x-12$;
\item $2x^2-3x-5$;
\item $a^3-2a^2-5a+6$;
\item $2x^2-3xy-2y^2$;
\item $(x^2-x-6)-(x^2+3x-4)=-4x-2$.
\end{enumerate}
\end{zadaniReseniBox}
\blok{Přechod k součinům více členů}
\begin{prikladReseniBox}[title={Příklad 3: trojčlen krát trojčlen}]
Uprav $(x^2+x+1)(x^2-x+1)$.
\tcblower
\begin{align*}
&(x^2+x+1)(x^2-x+1)\\
&=x^4-x^3+x^2+x^3-x^2+x+x^2-x+1\\
&=x^4+x^2+1.
\end{align*}
V devíti součinech je třeba hlídat i členy, které se vyruší.
\end{prikladReseniBox}
\blok{Volitelná rezerva}
\begin{zadaniReseniBox}[title={R1: urči chybu}]
Žák napsal $(x-2)(x+3)=x^2+3x-6$. Vysvětli chybu a výsledek oprav.
\tcblower
Vynechal součin $(-2)\cdot x=-2x$. Správně $x^2+3x-2x-6=x^2+x-6$.
\end{zadaniReseniBox}
\blok{Domácí procvičení}
\poznamka{Úlohy a)--f) jsou běžné procvičení, g)--i) náročnější volitelná část. V každé úloze požaduj mezikrok roznásobení, nejen výsledek.}
\begin{zadaniReseniBox}[title={D1: procvičení násobení}]
\begin{enumerate}[label=\alph*)]
\item $(2x+1)(x-3)$;
\item $(a-2)(a^2-3a+4)$;
\item $(x+2y)(3x-y)$;
\item $(x+1)(x-2)-(x-1)(x+2)$;
\item $(3x-4)(2x+1)$;
\item $(2a-b)(a^2+ab-b^2)$;
\item $(x^2-2x+1)(x+3)$;
\item $(x^2+x-1)(x^2-x+2)$;
\item $(2x-1)(x+2)-(x-3)(x+4)$.
\end{enumerate}
\tcblower
\begin{enumerate}[label=\alph*)]
\item $2x^2-5x-3$;
\item $a^3-5a^2+10a-8$;
\item $3x^2+5xy-2y^2$;
\item $-2x$;
\item $6x^2-5x-4$;
\item $2a^3+a^2b-3ab^2+b^3$;
\item $x^3+x^2-5x+3$;
\item $x^4+3x-2$;
\item $x^2+2x+10$.
\end{enumerate}
\end{zadaniReseniBox}
\end{document}
